\section{Nuestro método: \model}
Nuestro objetivo es diseñar un modelo $\Theta^\prime$ que preserve el rendimiento del Transformer $\Theta$ siendo al mismo tiempo más eficiente. Abordamos este problema aproximando las capas causal transformer en los video LMMs preentrenados.

\subsection{Actualización de tokens de texto mediante Cross-Attentions}
\label{sec:cross_attention}
La idea clave de nuestro método es dividir la costosa operación self-attention sobre toda la secuencia de video y text tokens en dos componentes más eficientes. Como los video tokens suelen dominar la secuencia mientras que los text tokens siguen siendo pocos, mantenemos el mecanismo self-attention exclusivamente para los tokens de texto y lo eliminamos para los video tokens. En su lugar, añadimos capas cross-attention que usan los tokens de texto como queries y los video tokens como keys y values. Como se muestra en la Figura~\ref{fig:main_attention}, este diseño actualiza los tokens de texto durante las capas de atención de nuestro modelo de la siguiente manera:
\vspace{-0.5em}
\begin{equation}
\label{equ:our_attention}
\begin{aligned}
    o_{t_j}= & (1-\alpha)(\sigma(\frac{q_{t_j} \mathbf{K}_v^\top}{\sqrt{d}})\mathbf{V}_v)\mathbf{W}^{c}_{o} & \text{\small(Cross-Attention)}\\ 
    & +\alpha(\sigma(\frac{q_{t_j} \mathbf{K}_{[t_1:t_j]}^\top}{\sqrt{d}})\mathbf{V}_{[t_1:t_j]}])\mathbf{W}^{s}_{o}, & \text{\small(Self-Attention)}
\end{aligned}
\end{equation}
donde $\mathbf{W}_o^c$ y $\mathbf{W}_o^s$ son las matrices de proyección de salida para las capas cross-attention y self-attention. $\alpha\in [0, 1]$ es un parámetro aprendible de ponderación que equilibra las salidas de cross-attention y self-attention. Suponiendo un total de $M$ video tokens y $N$ text tokens, las operaciones self-attention en los video LMMs basados en Transformer tienen una complejidad computacional de (pre-filling) $O(d(M+N)^2)$. Nuestro diseño self$+$cross-attention tiene complejidad $O(dN^2)$ (self-attention) y $O(dMN)$ (cross-attention), reduciendo efectivamente la complejidad de pre-filling a $O(dN^2+dMN)\sim O(dMN)$, dado que $M\gg N \approx d$.

Nuestra formulación de atención puede verse como una aproximación de la operación causal self-attention (Ecuación~\ref{equ:causal_self_attention_text}), ya que cada text token $q_{t_j}$ todavía puede atender a todos los video tokens y a los primeros $j$ text tokens. Como los tokens de video y texto comparten la misma dimensión de canal, las matrices de proyección en la capa cross-attention tendrán dimensiones idénticas a las de las capas self-attention. Esto nos permite considerar dos decisiones de diseño: podemos inicializar aleatoriamente los pesos en las capas cross-attention, o podemos inicializar las capas cross-attention usando los pesos de las capas self-attention (es decir, $\mathbf{W}_o^c=\mathbf{W}_o^s$, y lo mismo aplica a las matrices de proyección query, key y value $\mathbf{W}_q^c,\mathbf{W}_k^c,\mathbf{W}_v^c$). \textit{Añadimos estas dos decisiones al espacio de diseño de \model.}

\subsection{Actualización de video tokens con bloques Mamba}
Aunque el diseño self$+$cross-attention reduce significativamente la complejidad del modelo, depender únicamente de capas cross-attention puede comprometer la capacidad representacional del modelo. En concreto, en lugar de actualizar los video tokens mediante capas self-attention y MLP (es decir, a través del bloque causal transformer, \textit{c.f.} Sección~\ref{sec:autoregressive_lmm}), los video tokens permanecen sin cambios después de las capas cross-attention. Para abordar esta limitación, buscamos una arquitectura eficiente que aproxime los efectos de los bloques Transformer. Motivados por el éxito de las arquitecturas Mamba \cite{gu2023mamba, dao2024transformers} en image y video modelling \cite{zhu2024vision, liu2024vmamba, li2024videomamba}, proponemos emplear bloques Mamba para procesar eficazmente los video tokens. Como se muestra en la Figura~\ref{fig:main_attention}, las actualizaciones de video tokens pueden formularse como:
\begin{equation}
\label{equ:mamba}
    o_{v_i} = \mathrm{Mamba}(\mathrm{LN}(v_i), \mathbf{h}_{v_{i-1}}, \overline{\mathbf{A}}, \overline{\mathbf{B}}, \mathbf{C}),
\end{equation}
donde $\mathrm{LN}(\cdot)$ es el operador layer normalization y $\mathbf{h}_{v_{i-1}}$ es el contexto (hidden states) en el token $v_{i-1}$. Como modelo recurrente, Mamba opera de forma similar a las capas causal self-attention, permitiendo que cada video token se actualice a sí mismo a partir de todos los tokens anteriores. De manera crucial, Mamba reduce la complejidad de actualizar video tokens desde $O(dM^2)$ en causal self-attentions hasta $O(d^2M)$. Al combinar los bloques Mamba con el diseño self$+$cross-attention, nuestro modelo alcanza una complejidad total de $O(dN^2+dMN+d^2M)\sim O(dMN+d^2M)$, sustancialmente menor que la complejidad cuadrática $O(d(M+N)^2)$ en los video LMMs basados en Transformer. En la Sección~\ref{sec:training_paradigm}, detallamos además cómo se entrenan las capas Mamba para aproximar la Ecuación~\ref{equ:causal_self_attention_video}.

\textit{We consider employing Mamba or Mamba-2 blocks in Equation~\ref{equ:mamba} in the \model design space.}

\subsection{Paradigma de entrenamiento}
\label{sec:training_paradigm}
Utilizamos una estrategia de entrenamiento en dos etapas para optimizar \model, que incluye una etapa de pretraining y una etapa de instruction-tuning. En la etapa de pretraining, el modelo se inicializa con pesos preentrenados de un LMM basado en Transformer en todos los módulos excepto en las nuevas capas cross-attention y Mamba. Congelamos los componentes preentrenados y entrenamos solo las capas nuevas usando datos de image captioning para restaurar las capacidades de comprensión visual del modelo. Exploramos dos tipos de objetivos de entrenamiento durante el pretraining: (1) la language modelling loss estándar $\mathcal{L}_{\mathrm{LM}}=-\frac{1}{T}\sum_{t=1}^{T}\log p(x_t|x_{<t})$, que es la pérdida de cross-entropy para la predicción del siguiente token ($x_t$); (2) la distillation loss $\mathcal{L}_{\mathrm{Distill}}$, donde extraemos los 100 logits superiores $\mathcal{P}_{\Theta}$ con los valores más altos del modelo basado en Transformer $\Theta$ y calculamos una pérdida KL-divergence con las salidas logit $\mathcal{P}_{\Theta^\prime}$ de nuestro modelo $\Theta^\prime$ en los mismos índices, de modo que $\mathcal{L}_{\mathrm{Distill}}=D_{KL}(\mathcal{P}_{\Theta}||\mathcal{P}_{\Theta^\prime})$
($D_{KL}(\cdot||\cdot)$ denota KL-divergence). La pérdida final se formula como $\mathcal{L}=\mathcal{L}_{\mathrm{LM}} + \lambda\mathcal{L}_{\mathrm{Distill}}$, donde usamos $\lambda$ para equilibrar los pesos de ambas pérdidas. \textit{Incluimos $\lambda=0, 0.001, 0.1, 0.5, 1, 2$ en el espacio de diseño de \model.}

En la etapa de instruction-tuning, aprovechamos tanto datos de image como de video instruction-following para hacer full finetune de \model, mejorando así su capacidad de instruction-following. Cuando las limitaciones de memoria GPU impiden el full finetuning, congelamos el vision encoder y ajustamos únicamente el decoder LMM. Durante la etapa de instruction-tuning empleamos solo la language modelling loss para asegurar que el teacher model no restrinja el rendimiento del student.

\textit{El espacio de diseño completo de \model incluye si inicializar los pesos cross-attention a partir de self-attention; si usar bloques Mamba o Mamba-2 para actualizar video tokens; y la elección de $\lambda$ en la etapa de pretraining.}
